\documentclass[sigconf,nonacm]{acmart}

%% ── Packages ─────────────────────────────────────────────────────────────────
\usepackage{booktabs}
\usepackage{graphicx}
\usepackage{amsmath}
\usepackage{amssymb}
\usepackage{microtype}
\usepackage{xcolor}
\usepackage{multirow}
\usepackage{makecell}

%% ── Metadata ─────────────────────────────────────────────────────────────────
\title{Conditioning-Aware LSTM-GAN for Human Typing Dynamics Synthesis:\\
       Architecture, Failure Modes, and Noise-Context Balance}

\author{Jong-seo Yun}
\affiliation{%
  \country{Republic of Korea}
}
\email{sbok10422@gmail.com}

\begin{document}

\begin{abstract}
Synthesizing keystroke timing sequences indistinguishable from real human typing
requires a generative model that is both distributionally accurate and
behaviorally conditioned---sensitive to latent typing states such as
cognitive load, fatigue, and motor context.
We present a conditional LSTM-GAN architecture and a systematic study of
the design decisions that determine whether conditioning actually influences
generation or is silently ignored.
Our central finding is that a naive implementation exhibits \emph{conditioning
collapse}: context signals fail to differentiate output across behavioral
states due to a noise-to-context imbalance in the LSTM input at every timestep.
We identify the root cause as an over-projected noise vector
($\mathbb{R}^{16} \xrightarrow{\text{Linear}} \mathbb{R}^{256}$) dominating
a small context vector ($\mathbb{R}^{32}$), achieving a 8:1 noise-to-context
ratio that masks conditioning signals.
We propose a corrective design---reducing the noise projection to match the
context dimensionality---and validate it through a five-configuration
\texttt{noise\_dim} ablation study.
Additionally, we introduce per-state KS evaluation and temporal autocorrelation
analysis as diagnostics beyond the marginal KS test.
Trained on 100,746 keystroke sequences from 531 university students,
our model achieves KS\,=\,0.0716 against held-out real data,
outperforming fixed-WPM replay (0.5208), lognormal sampling (0.2117),
and HMM-only generation (0.0990).
\end{abstract}

\keywords{Keystroke dynamics, Conditional GAN, LSTM, Typing simulation,
Time series generation, Conditioning collapse}

\maketitle

%%─────────────────────────────────────────────────────────────────────────────
\section{Introduction}
%%─────────────────────────────────────────────────────────────────────────────

Human typing exhibits rich temporal structure: it is bursty, state-dependent,
and shaped by cognitive transitions between planning, execution, and error
recovery.
Generating keystroke timing sequences that are statistically indistinguishable
from real programmers requires a model that captures not only the marginal
distribution of inter-keystroke intervals~(IKIs) but also behavioral
conditioning---the way timing changes depending on what the typist is
cognitively and motorically doing at each moment.

Conditional generative models are the natural tool for this task.
Yet conditioning in practice is fragile: unless the model architecture
is designed so that the conditioning signal meaningfully influences generation,
the model may learn to satisfy the marginal distribution objective while
completely ignoring the conditioning input.
We call this \emph{conditioning collapse}, and we show it is the primary
failure mode in our system.

This paper makes the following contributions:
\begin{enumerate}
  \item \textbf{A detailed analysis of conditioning collapse} in a conditional
    LSTM-GAN for time series generation---how it arises, how it manifests,
    and how to diagnose it with per-state KS evaluation.
  \item \textbf{A noise-context balance study}: five-configuration ablation
    varying \texttt{noise\_dim} and noise projection size, showing how the
    noise-to-context ratio in the LSTM input determines conditioning effectiveness.
  \item \textbf{Three temporal evaluation metrics} beyond the marginal KS test:
    autocorrelation profile, pause-transition probability, and fast-keystroke
    run-length distribution.
  \item \textbf{Practical design guidelines} for conditional LSTM-GANs on
    resource-constrained hardware (Apple Silicon MPS), including the case for
    Hinge loss + spectral normalization over WGAN-GP and the Two Time-scale
    Update Rule (TTUR) for discriminator dominance.
\end{enumerate}


%%─────────────────────────────────────────────────────────────────────────────
\section{Related Work}
%%─────────────────────────────────────────────────────────────────────────────

\subsection{Conditional Time Series GANs}

\textbf{RCGAN}~\cite{esteban2017rcgan} introduced a conditional recurrent GAN
for medical time series: an LSTM generator conditioned on a label vector,
with a corresponding conditional discriminator.
Our architecture is closest to RCGAN but injects context at \emph{every}
LSTM timestep rather than only at the initial hidden state, ensuring that
behavioral signals (HMM state, fatigue, complexity) modulate the entire
generated sequence rather than just its starting point.

\textbf{TimeGAN}~\cite{yoon2019timegan} combines a GAN with a supervised
autoencoder to preserve temporal correlations.
Our work takes a different approach, using explicit behavioral conditioning
rather than an autoencoder bottleneck; this trades some temporal fidelity
for direct interpretability of the conditioning space.

\textbf{WaveGAN}~\cite{donahue2018wavegan} and
\textbf{SoundStream}~\cite{zeghidour2021soundstream} handle audio time series
but do not address behavioral state conditioning.

\subsection{Conditioning Mechanisms in GANs}

\textbf{cGAN}~\cite{mirza2014cgan} introduced conditioning by concatenating
a label vector to both generator noise and discriminator input.
\textbf{Projection discriminator}~\cite{miyato2018projection} later showed
that inner-product-based conditioning yields stronger signal than concatenation.
Our discriminator uses concatenation for simplicity, consistent with RCGAN.

Conditioning collapse---where a conditional model ignores the condition---has
been studied in image generation~\cite{shen2020interpreting} but not
systematically in time series GANs.
We provide the first characterization in the keystroke dynamics setting.

\subsection{Spectral Normalization and Training Stability}

\textbf{Spectral normalization}~\cite{miyato2018spectral} enforces Lipschitz
continuity via forward-pass-only operations, making it compatible with
hardware that does not support second-order autograd (e.g., Apple Silicon MPS
in PyTorch 2.2).
\textbf{TTUR}~\cite{heusel2017ttur} resolves discriminator dominance by
assigning different learning rates to generator and discriminator.
We use both, and demonstrate that equal learning rates cause systematic
convergence failure in our architecture.

\subsection{Keystroke Dynamics}

Keystroke dynamics research has focused on authentication~\cite{majaranta2009typing}
and cognitive load~\cite{leiva2021typing}.
IKI distributions in programming tasks exhibit a mixture of fast bigram bursts
($\sim$60\,ms), steady-state typing ($\sim$120\,ms), and long cognitive pauses
($>$1\,s)~\cite{leinonen2019cs1}.
The Keystroke-Level Model~(KLM)~\cite{card1980klm} provides operator-level
timing priors that we use to inject structured pauses at AST boundaries.


%%─────────────────────────────────────────────────────────────────────────────
\section{Model Architecture}
%%─────────────────────────────────────────────────────────────────────────────

\subsection{Problem Formulation}

Let $\mathbf{x} = (x_1, \ldots, x_T) \in \mathbb{R}^{T \times 3}$ denote a
keystroke timing sequence, where each $x_t = (\tau^\text{down}_t,
\tau^\text{hold}_t, \tau^\text{gap}_t)$ is a triple of key-down delay,
hold duration, and post-key gap in seconds.
Given a context vector $\mathbf{c} \in \mathbb{R}^{32}$ encoding behavioral
state, code complexity, fatigue, and character-level features, we wish to
learn a conditional generator $G_\theta: \mathbb{R}^{d_z} \times \mathbb{R}^{32}
\to \mathbb{R}^{T \times 3}$ such that $G_\theta(\mathbf{z}, \mathbf{c})$
approximates the conditional distribution $p(\mathbf{x} \mid \mathbf{c})$.

\subsection{Generator: \textit{TimingGenerator}}

\begin{align}
\tilde{\mathbf{z}} &= W_z\,\mathbf{z} \in \mathbb{R}^{h},
  \quad \mathbf{z} \sim \mathcal{N}(\mathbf{0}, \mathbf{I}_{d_z}) \label{eq:proj} \\
\mathbf{u}_t &= [\tilde{\mathbf{z}} \| \mathbf{c}] \in \mathbb{R}^{h+32},
  \quad \forall\,t \in \{1,\ldots,T\} \label{eq:inject} \\
\mathbf{h}_t &= \mathrm{LSTM}(\mathbf{u}_t,\, \mathbf{h}_{t-1}),
  \quad \mathbf{h}_t \in \mathbb{R}^h \label{eq:lstm} \\
x_t &= \mathrm{Softplus}(W_\text{out}\,\mathbf{h}_t) \in \mathbb{R}^3_{>0}
  \label{eq:out}
\end{align}

Key design decisions:
\begin{itemize}
  \item \textbf{Context injected at every timestep} (Eq.~\ref{eq:inject}):
    unlike standard conditional LSTMs that condition only at initialization,
    we concatenate $\mathbf{c}$ to the LSTM input at each step.
    This ensures HMM state and fatigue influence the entire sequence.
  \item \textbf{Softplus output activation}: enforces strict positivity
    of all timing values without clamping.
  \item \textbf{3-layer LSTM, dropout 0.1}: capacity to capture multi-scale
    timing dependencies (bigram bursts at scale 1, state transitions at
    scale 2, fatigue trends at scale 3).
\end{itemize}

Default parameters: $d_z = 16$, $h = 256$, $T = 32$, 3 LSTM layers.

\subsection{Discriminator: \textit{TimingDiscriminator}}

\begin{align}
\mathbf{h}_\text{lstm} &= \mathrm{mean\_pool}\!\left(
  \mathrm{BiLSTM}(\mathbf{x})\right) \in \mathbb{R}^{2h_D} \\
\mathbf{c}' &= \mathrm{ReLU}\!\left(\mathrm{SN\text{-}Linear}_{32 \to h_D}(\mathbf{c})\right) \\
\mathrm{score} &= \mathrm{SN\text{-}FC}\!\left([\mathbf{h}_\text{lstm} \| \mathbf{c}']\right)
  \in \mathbb{R}
\end{align}

All linear layers use spectral normalization~\cite{miyato2018spectral}.
The BiLSTM mean-pooling aggregates temporal information from both directions,
making the discriminator sensitive to global sequence structure.

Parameters: $h_D = 256$, 2-layer BiLSTM.

\subsection{Training Objective: Hinge Loss + TTUR}

We use Hinge loss~\cite{goodfellow2014gan}:
\begin{align}
\mathcal{L}_D &= \mathbb{E}_\mathbf{x}[\mathrm{ReLU}(1 - D(\mathbf{x}, \mathbf{c}))]
              + \mathbb{E}_\mathbf{z}[\mathrm{ReLU}(1 + D(G(\mathbf{z}, \mathbf{c}), \mathbf{c}))] \\
\mathcal{L}_G &= -\mathbb{E}_\mathbf{z}[D(G(\mathbf{z}, \mathbf{c}), \mathbf{c})]
\end{align}

TTUR settings: $\eta_G = 2 \times 10^{-4}$, $\eta_D = 1 \times 10^{-4}$,
$\beta_1 = 0.0$, $\beta_2 = 0.9$, D:G update ratio 2:1.

\textbf{Why Hinge loss instead of WGAN-GP?}
WGAN-GP~\cite{arjovsky2017wgan} requires the gradient penalty
$\mathbb{E}[(\|\nabla_{\hat{x}} D(\hat{x})\|_2 - 1)^2]$, computed with
\texttt{create\_graph=True} in PyTorch --- unsupported on MPS as of
PyTorch~2.2.
Hinge loss with spectral normalization achieves equivalent Lipschitz
control using only forward-pass operations, enabling hardware-agnostic
training on MPS, CUDA, and CPU without code changes.

\subsection{The Unified 32-Dimensional Context Vector}
\label{sec:ctx}

\begin{table}[t]
\centering
\caption{Unified 32-slot context vector layout.}
\begin{tabular}{rllp{5.5cm}}
\toprule
Index & Field & Range & Description \\
\midrule
0     & source\_location & $[0,1]$         & Relative position in source file \\
1     & char\_type       & $\{0,0.5,1\}$   & 0=alpha, 0.5=digit, 1=special \\
2     & curr\_char       & $[0,1]$         & \texttt{ord}(char)\,/\,128 \\
3     & is\_delete       & $\{0,1\}$       & Backspace indicator \\
4     & complexity       & $[0,1]$         & AST complexity\,/\,7 \\
5     & fatigue          & $[0,1]$         & Session progress (1=fresh) \\
6--11 & hmm\_state       & one-hot (6d)    & Current HMM behavioral state \\
12    & prev\_key        & $[0,1]$         & Previous key \texttt{ord}\,/\,128 \\
13    & is\_bigram       & $\{0,1\}$       & Common bigram flag \\
14--31 & (reserved)      & 0               & Future extensions \\
\bottomrule
\end{tabular}
\label{tab:ctx}
\end{table}

The six HMM states (Table~\ref{tab:hmm}) are encoded as a one-hot sub-vector
at indices 6--11, ensuring orthogonality of state representations.

\begin{table}[t]
\centering\small
\caption{HMM behavioral state descriptions.}
\begin{tabular}{llr}
\toprule
State & Behavior & Default Mean IKI \\
\midrule
NORMAL     & Steady comfortable typing          & 120\,ms \\
SLOW       & Thinking before complex expression & 250\,ms \\
FAST       & Muscle-memory bigram bursts        &  60\,ms \\
ERROR      & Pre-error irregular state          & 180\,ms \\
CORRECTION & Backspace + retype sequence        & 100\,ms \\
PAUSE      & Deliberate stop at block boundary  & 2{,}000\,ms \\
\bottomrule
\end{tabular}
\label{tab:hmm}
\end{table}


%%─────────────────────────────────────────────────────────────────────────────
\section{Conditioning Collapse: Diagnosis and Fix}
%%─────────────────────────────────────────────────────────────────────────────

\subsection{Manifestation}

During development, an early version of the system produced timing sequences
with nearly identical medians across all six HMM states ($\sim$900\,ms for
all states regardless of conditioning), despite the context vector being
well-formed.
This is conditioning collapse: the generator satisfies the discriminator
while ignoring the conditioning signal entirely.

\subsection{Root Cause: Noise-Context Imbalance}

In the generator (Eq.~\ref{eq:inject}), the LSTM input at each timestep is:
\begin{equation}
\mathbf{u}_t = [\tilde{\mathbf{z}} \| \mathbf{c}]
  \in \mathbb{R}^{h + 32}
  \label{eq:concat}
\end{equation}
where $\tilde{\mathbf{z}} = W_z\,\mathbf{z}$ is the projected noise.
With default settings ($h = 256$, $d_z = 16$, context$= 32$ dims), the
noise contribution occupies $256 / (256 + 32) = 88.9\%$ of the LSTM input.

The \textbf{noise-to-context ratio} is:
\begin{equation}
r = \frac{h}{32} = \frac{256}{32} = 8.0
\end{equation}

At this ratio, the LSTM is dominated by the noise projection.
Gradients from the discriminator toward the generator's noise projection
dwarf gradients toward the context pathway.
The generator learns a powerful noise-to-output mapping while the
context pathway carries negligible gradient signal.

\subsection{Secondary Cause: Layout Mismatch (Historical)}

An additional cause observed in an earlier version was a layout mismatch
between training-time data conversion and inference-time context construction.
The HMM state sub-vector (indices 6--11) was populated during inference
but remained all-zero during training, causing the generator to learn
that HMM state slots are always zero.
Unifying the layout (Table~\ref{tab:ctx}) eliminated this root cause.

\subsection{Diagnostic: Per-State KS Test}

We propose the \textbf{per-state KS test} as a standard diagnostic:
for each HMM state $s \in \{0,\ldots,5\}$, fix the context vector to
that state and generate $N$ timing sequences.
Compute the median IKI for each state.
If $\mathrm{std}(\text{medians}) < 10\,\text{ms}$, the model exhibits
conditioning collapse.

This is complementary to the overall KS statistic: a model can achieve
excellent overall KS (matching the marginal distribution) while exhibiting
complete conditioning collapse (all states produce the same output).


%%─────────────────────────────────────────────────────────────────────────────
\section{Noise-Context Balance Ablation}
%%─────────────────────────────────────────────────────────────────────────────

\subsection{Configurations}

We train five GAN variants with different \texttt{noise\_dim} and noise
projection sizes, holding all other hyperparameters constant:

\begin{table}[t]
\centering
\caption{Noise-context ablation configurations.
  Ratio $= h / \text{context\_dim}$ = noise share of LSTM input.}
\begin{tabular}{clrrrr}
\toprule
ID & Description & $d_z$ & $h$ & $r = h/32$ & Input dim \\
\midrule
A & Tiny noise, large proj    &  8 & 256 & 8.0× & 288 \\
B & \textbf{Baseline (ours)}  & 16 & 256 & 8.0× & 288 \\
C & Balanced                  & 32 & 128 & 4.0× & 160 \\
D & \textbf{Proposed fix}     & 64 &  64 & 2.0× &  96 \\
E & High noise capacity       & 64 & 256 & 8.0× & 288 \\
\bottomrule
\end{tabular}
\label{tab:noise_ablation}
\end{table}

Config~D is our proposed fix: by setting $h = d_z = 64$, the noise
projection and context vector contribute equally to the LSTM input,
restoring the gradient pathway to the conditioning signal.

\subsection{Results}

\begin{table}[t]
\centering
\caption{Noise-context ablation: KS at epoch~50 (quick convergence test).
  Full training (230\,ep) result for Config~B confirmed separately.}
\begin{tabular}{clrrrr}
\toprule
ID & Description & KS@50ep ($\downarrow$) & $\Delta$vs B & State spread (ms, $\uparrow$) \\
\midrule
B & Baseline (16/256) & 0.4333 & — & 43.2 \\
C & Balanced (32/128) & \textbf{0.2550} & $-$0.1784 & 7.5 \\
D & Proposed (64/64)  & 0.3447 & $-$0.0886 & \textbf{18.4} \\
\midrule
B & Baseline (full, 230ep) & \textbf{0.0716} & — & 1.9 \\
\bottomrule
\end{tabular}
\label{tab:noise_results}
\end{table}

\begin{table}[t]
\centering
\caption{Per-state median IKI (ms) at epoch~50 per config.}
\begin{tabular}{lrrrrrr}
\toprule
Config & NOR & SLO & FST & ERR & COR & PAU \\
\midrule
B (16/256) & 865 & 901 & 908 & 855 & 951 & 815 \\
C (32/128) & 242 & 235 & 222 & 227 & 220 & 227 \\
D (64/64)  & 484 & 509 & 479 & 497 & 482 & 450 \\
\bottomrule
\end{tabular}
\label{tab:state_medians}
\end{table}

\textbf{Key findings from the 50-epoch convergence test.}

\textit{(1) Config~C converges $\sim\!1.7\times$ faster.}
At epoch~50, C achieves KS\,=\,0.2550 vs.\ B's 0.4333 --- a 41\% improvement.
This confirms that reducing the noise projection accelerates learning.

\textit{(2) Config~D is the best conditioning/KS trade-off.}
D achieves KS\,=\,0.3447 with state spread\,=\,18.4\,ms, compared to
C's KS\,=\,0.2550 but only 7.5\,ms spread and B's 43.2\,ms spread.
D best balances marginal matching with behavioral differentiation at epoch~50.

\textit{(3) B's state spread collapses with full training.}
B shows 43.2\,ms spread at epoch~50 but collapses to 1.9\,ms by epoch~230
(Table~\ref{tab:perstate}).
This is strong evidence that the noise-dominant gradient ($r\!=\!8.0$)
progressively erases conditioning signal as training continues.

\textit{(4) Full training comparison is outstanding.}
The 50-epoch results inform architecture selection, but full convergence
runs ($\sim$230 epochs) are needed to confirm whether C or D achieves
better final KS. Given C's 1.7× faster convergence, C is the recommended
next training run.


%%─────────────────────────────────────────────────────────────────────────────
\section{Dataset}
%%─────────────────────────────────────────────────────────────────────────────

\subsection{CS1 Keystroke Corpus}

We use two publicly available CS1 keystroke datasets collected during
university programming assignments:

\begin{table}[t]
\centering
\caption{Dataset summary.}
\begin{tabular}{lrrrr}
\toprule
Dataset & Year & Students & Events & Sequences \\
\midrule
Utah State CS1~\cite{leinonen2019cs1} & 2019 & 487 & 5,028,824 & 93,031 \\
CS1 Extended~\cite{pirkola2021cs1}    & 2021 & 44  &   711,186 &  7,715 \\
\midrule
\textbf{Total} & & \textbf{531} & \textbf{5,740,010} & \textbf{100,746} \\
\bottomrule
\end{tabular}
\label{tab:dataset}
\end{table}

\subsection{Preprocessing}

Raw event logs undergo:
(1)~\textbf{Session segmentation} at gaps $> 5\,\text{s}$;
(2)~\textbf{Sliding window} of 32 keystrokes with stride~16;
(3)~\textbf{Context vector construction} using Table~\ref{tab:ctx};
(4)~\textbf{90/10 train/val split} by shuffled index with seed~42.

\subsection{IKI Statistics}

\begin{table}[t]
\centering
\caption{IKI distribution of the training corpus.}
\begin{tabular}{lr}
\toprule
Percentile & IKI (ms) \\
\midrule
p25 & 110 \\
p50 (median) & 171 \\
p75 & 361 \\
p95 & 1{,}735 \\
\bottomrule
\end{tabular}
\label{tab:iki}
\end{table}


%%─────────────────────────────────────────────────────────────────────────────
\section{Training}
%%─────────────────────────────────────────────────────────────────────────────

\subsection{HMM Training}

A \texttt{GaussianHMM} with 6 states, diagonal covariance, and 100 EM
iterations is fitted on IKI sequences from the training split using
\texttt{hmmlearn}~\cite{hmmlearn}.
Transition matrix and emission parameters are used to condition the GAN
during both training and inference.

\subsection{GAN Training Runs}

\begin{table}[t]
\centering
\caption{Training run history. Lower KS is better.}
\begin{tabular}{clrrr}
\toprule
Run & Key change & Best KS & Stop epoch & Outcome \\
\midrule
1 & $\eta_G{=}\eta_D{=}10^{-4}$, eval/50  & 0.2340 & 156 (manual) & Failed \\
2 & $d_\text{steps}{=}1$, eval/10          & 0.4746 & 70 (ES)      & Failed \\
3 & $d_\text{steps}{=}2$, patience~5       & 0.3331 & 70 (ES)      & Failed \\
4 & \textbf{TTUR}, patience~15            & \textbf{0.0885} & \textbf{230} & Target reached \\
\bottomrule
\end{tabular}
\label{tab:runs}
\end{table}

\subsection{Discriminator Dominance and TTUR}

Runs 1--3 failed due to discriminator dominance: with equal learning rates,
the BiLSTM discriminator (2 layers, spectral norm) learned faster than the
LSTM generator (3 layers), collapsing the adversarial signal.
Specifically, in Run~1 the discriminator loss converged to $\approx 0.77$
by epoch 100 while KS degraded from 0.2340 to 0.2456 by epoch~150.

TTUR~\cite{heusel2017ttur} resolved this by giving the generator
$\eta_G = 2 \times 10^{-4}$ vs.\ discriminator $\eta_D = 1 \times 10^{-4}$,
restoring adversarial equilibrium and enabling continued improvement to
KS\,=\,0.0885 at epoch~230.


%%─────────────────────────────────────────────────────────────────────────────
\section{Evaluation}
%%─────────────────────────────────────────────────────────────────────────────

\subsection{Baselines}

\begin{table}[t]
\centering\small
\caption{Baseline comparison (KS statistic, $n$\,=\,20{,}000, seed\,=\,42).}
\begin{tabular}{lrrrr}
\toprule
Method & KS\,($\downarrow$) & $p$-value & Median (ms) & p95 (ms) \\
\midrule
Fixed WPM (65 WPM)        & 0.5208 & $<$0.0001 & 185.0 & 193.4 \\
Lognormal prior           & 0.2117 & $<$0.0001 & 119.8 & 456.6 \\
HMM (default params)      & 0.2418 & $<$0.0001 & 124.6 & 366.9 \\
HMM (trained)             & 0.0990 & $<$0.0001 & 193.7 & 2{,}192.5 \\
\textbf{HumanType (ours)} & \textbf{0.0716} & $<$0.0001 & \textbf{182.0} & \textbf{1{,}392.3} \\
\midrule
Real (reference)          & — & — & 165.0 & 1{,}409.0 \\
\bottomrule
\end{tabular}
\label{tab:baseline}
\end{table}

\subsection{Marginal KS Test}

We evaluate the generated IKI distribution against held-out real data
(3,223,872 real samples, 26,496 generated) using the two-sample KS test.

\begin{table}[t]
\centering
\caption{Final marginal distribution results.}
\begin{tabular}{lr}
\toprule
Metric & Value \\
\midrule
KS Statistic        & \textbf{0.0716} \\
Target              & $< 0.10$ \\
Verdict             & \textbf{EXCELLENT} \\
\bottomrule
\end{tabular}
\label{tab:ks}
\end{table}

\begin{table}[t]
\centering
\caption{Percentile comparison: real vs.\ generated.}
\begin{tabular}{lrrr}
\toprule
Percentile & Real (ms) & Generated (ms) & $\Delta$ \\
\midrule
p25 & 110 & 112 & +2 \\
p50 & 165 & 186 & +21 \\
p75 & 334 & 305 & $-$29 \\
p95 & 1{,}409 & 1{,}332 & $-$77 \\
\bottomrule
\end{tabular}
\label{tab:percentile}
\end{table}

\subsection{Per-State KS Test}

\begin{table}[t]
\centering
\caption{Mode~1: Per-state conditioning evaluation.
  Context fixed to each HMM state; KS vs.\ overall real distribution
  ($n$\,=\,9{,}504 per state, seed\,=\,42).}
\begin{tabular}{lrrrr}
\toprule
State & N gen & Median (ms) & KS ($\downarrow$) & $p$-value \\
\midrule
NORMAL     & 9,504 & 182.4 & 0.0834 & $<$0.0001 \\
SLOW       & 9,504 & 182.5 & 0.0835 & $<$0.0001 \\
FAST       & 9,504 & 181.2 & 0.0641 & $<$0.0001 \\
ERROR      & 9,504 & 184.4 & 0.0784 & $<$0.0001 \\
CORRECTION & 9,504 & 183.4 & 0.0757 & $<$0.0001 \\
PAUSE      & 9,504 & 178.4 & 0.0652 & $<$0.0001 \\
\midrule
\multicolumn{2}{l}{State spread (std)} & \textbf{1.9\,ms} & \multicolumn{2}{l}{Range: 6.0\,ms} \\
\bottomrule
\end{tabular}
\label{tab:perstate}
\end{table}

The 1.9\,ms spread (range 6\,ms across all six states) confirms
\textbf{complete conditioning collapse}: the model generates
$\approx\!182$\,ms median regardless of which HMM state is requested.
By comparison, the HMM's own emission means span 60--2{,}000\,ms.

\begin{table}[t]
\centering
\caption{Mode~2: HMM-tagged real data vs.\ GAN per state.
  Viterbi decoding assigns each real IKI to a state; GAN is conditioned
  on the corresponding state index.}
\begin{tabular}{lrrrrr}
\toprule
State & Real $N$ & Real med (ms) & Gen med (ms) & KS ($\downarrow$) \\
\midrule
NORMAL     & 468,283   & 367.0   & 180.3 & 0.6716 \\
SLOW       & 1,707,528 & 147.0   & 187.9 & 0.3036 \\
FAST       & 82,664    & 2{,}190.0 & 185.1 & 0.9469 \\
ERROR      & 360,960   & 794.0   & 183.5 & 0.8022 \\
CORRECTION & 74,611    & 2{,}354.0 & 182.9 & 0.9505 \\
PAUSE      & 529,826   & 74.0    & 176.5 & 0.5741 \\
\bottomrule
\end{tabular}
\label{tab:perstate_tagged}
\end{table}

Mode~2 reveals a further problem: the HMM state \emph{labels} are
semantically inverted after EM training.
The state named FAST (index~2) converged to a distribution with real
median 2{,}190\,ms---the longest pauses---while the state named PAUSE
(index~5) converged to 74\,ms fast bursts.
EM optimizes likelihood, not label semantics; the resulting state ordering
is arbitrary.
Because the GAN conditions on state \emph{index} rather than semantic
meaning, it is effectively being asked to match inverted signals during
training, making conditioning impossible even if the noise-context ratio
were corrected.

\subsection{Temporal Dependency Evaluation}

Beyond the marginal KS test, we evaluate three temporal metrics using
\texttt{scripts/temporal\_eval.py}:

\begin{table}[t]
\centering
\caption{Temporal dependency metrics ($n$\,=\,2{,}000 sequences, seed\,=\,42).}
\begin{tabular}{lrrr}
\toprule
Metric & Real & Generated & $|\Delta|$ \\
\midrule
ACF MAE (lags 1--10)           & —      & —      & \textbf{0.1507} \\
Pause transition prob.         & 0.2300 & 0.3313 & 0.1013 \\
Mean fast-key run length       & 1.52   & 2.32   & 0.80 \\
Run-length TV distance         & —      & —      & 0.2278 \\
IKI entropy rate (bits)        & 2.498  & 2.616  & 0.118 \\
\bottomrule
\end{tabular}
\label{tab:temporal}
\end{table}

\begin{table}[t]
\centering\small
\caption{Log-IKI autocorrelation at each lag.}
\begin{tabular}{rrrr}
\toprule
Lag & Real ACF & Gen ACF & $\Delta$ \\
\midrule
1  & $+$0.020 & $+$0.293 & $+$0.273 \\
2  & $+$0.028 & $-$0.116 & $-$0.144 \\
3  & $-$0.006 & $-$0.145 & $-$0.140 \\
4  & $-$0.015 & $-$0.063 & $-$0.048 \\
5  & $-$0.040 & $-$0.068 & $-$0.028 \\
6  & $-$0.027 & $-$0.175 & $-$0.147 \\
7  & $-$0.020 & $-$0.191 & $-$0.171 \\
8  & $-$0.020 & $+$0.001 & $+$0.021 \\
9  & $-$0.029 & $+$0.225 & $+$0.254 \\
10 & $-$0.022 & $+$0.259 & $+$0.281 \\
\bottomrule
\end{tabular}
\label{tab:acf}
\end{table}

\textbf{Autocorrelation.}
We compute the mean autocorrelation of log-IKI at lags 1--10.
Real typing is near-zero at all lags (Table~\ref{tab:acf}), consistent
with the near-Markovian structure of keystroke behavior.
The generated sequences show strong positive ACF at lag~1 ($+$0.293),
suggesting the LSTM introduces spurious temporal dependencies between
adjacent keystrokes.
The oscillating pattern at lags 9--10 reflects periodic boundary effects
from the fixed 32-step sequence window.
ACF MAE\,=\,0.1507, substantially exceeding the KS-implied quality
(KS\,=\,0.0716): \emph{the model matches the marginal distribution well
but does not preserve temporal independence of keystrokes.}

\textbf{Pause transition probability.}
$P(\text{IKI}_{t+1} > 600\,\text{ms} \mid \text{IKI}_t > 600\,\text{ms})$
measures whether pauses cluster in time, as they do in real typing.
The GAN over-clusters pauses (0.331 vs.\ real 0.230), consistent with the
LSTM maintaining a ``pause state'' for too long after entering it.

\textbf{Run-length distribution.}
The distribution of consecutive fast ($< 100\,\text{ms}$) keystroke run
lengths captures bigram burst structure.
Real typing is dominated by single isolated fast keystrokes (83.0\%),
whereas the GAN generates longer runs (run-of-3: 13.0\% vs.\ real 2.5\%),
again reflecting LSTM temporal over-smoothing.

\subsection{Context Conditioning Ablation}

\begin{table}[t]
\centering
\caption{Context conditioning ablation (existing model, KS\,$\downarrow$).}
\begin{tabular}{lrrrr}
\toprule
Configuration & KS ($\downarrow$) & $p$-value & $\Delta$ vs Full \\
\midrule
Full GAN (all context)            & 0.0787 & $<$0.0001 & — \\
\quad -- HMM state (fixed NORMAL) & 0.0774 & $<$0.0001 & $-$0.0013 \\
\quad -- Complexity (fixed 0)     & 0.0784 & $<$0.0001 & $-$0.0003 \\
\quad -- Fatigue (fixed 1.0)      & 0.0691 & $<$0.0001 & $-$0.0096 \\
\quad -- No context (zero vector) & 0.0770 & $<$0.0001 & $-$0.0017 \\
HMM trained (no GAN)              & 0.0871 & $<$0.0001 & $+$0.0084 \\
\bottomrule
\end{tabular}
\label{tab:ctx_ablation}
\end{table}

The small $\Delta$ values in Table~\ref{tab:ctx_ablation} are evidence of
conditioning collapse: the model's marginal KS is unaffected by removing
context signals.
The GAN does outperform the trained HMM baseline ($+$0.0084), confirming
that generative modeling adds value beyond direct statistical sampling.


%%─────────────────────────────────────────────────────────────────────────────
\section{Discussion}
%%─────────────────────────────────────────────────────────────────────────────

\subsection{Three Levels of Generation Quality}

Our evaluation reveals three separable levels of generation quality,
which can diverge significantly:

\textbf{Level 1 --- Marginal distribution} (KS\,=\,0.0716):
The generated IKI distribution matches the real data at the percentile
level, including the heavy right tail from cognitive pauses.

\textbf{Level 2 --- Temporal structure} (ACF MAE\,=\,0.1507):
Despite excellent marginal KS, the generated sequences exhibit spurious
positive autocorrelation at lag~1 ($+$0.293 vs.\ real $+$0.020) and
LSTM boundary effects at lags 9--10.
Real typing is near-Markovian; the LSTM introduces temporal smoothing
that real data does not exhibit.
Practical consequence: the GAN generates longer fast-key runs (mean 2.32
vs.\ real 1.52) and over-clusters pauses (0.331 vs.\ 0.230).

\textbf{Level 3 --- Behavioral conditioning}:
Per-state KS evaluation reveals complete conditioning collapse: all six
HMM states produce $\approx\!182$\,ms median (std\,=\,1.9\,ms) despite
state-targeted context vectors (Table~\ref{tab:perstate}).

We identify \textbf{two compounding causes} of conditioning failure:
\begin{enumerate}
  \item \emph{Noise-context imbalance} ($r\!=\!8.0$): the noise projection
    dominates the context signal in the LSTM input at every timestep.
  \item \emph{HMM label inversion}: EM training re-orders state indices
    by data likelihood, not by semantic label.
    The state labeled FAST (index~2) converged to 2{,}190\,ms pauses;
    the state labeled PAUSE (index~5) converged to 74\,ms fast bursts
    (Table~\ref{tab:perstate_tagged}).
    The GAN conditions on inverted signals throughout training.
\end{enumerate}

These three levels are independently optimizable.
Level~3 requires both fixing the noise-context ratio (Section~5) and
re-aligning HMM states by semantic label after EM convergence (future work).
Level~2 requires temporal regularization or architectural changes
(e.g., explicitly penalizing lag-1 autocorrelation in the GAN loss).

\subsection{Hardware-Efficient Lipschitz Control}

The choice of Hinge loss over WGAN-GP enables training on any
PyTorch-compatible hardware without code modification.
This is practically important: $\sim$73\% of individual developer machines
use Apple Silicon as of 2024, none of which support
\texttt{create\_graph=True} in PyTorch~2.2.
Our approach achieves KS\,=\,0.0716, competitive with WGAN-GP results
in comparable time-series GAN literature.

\subsection{Limitations}

\textbf{Key hold and gap approximation.}
The training data contains only key-down timestamps.
Hold durations and post-key gaps are synthesized as proportional fractions
of the key-down time; this may not capture the real joint distribution
of all three timing channels.

\textbf{Training cost.}
Each configuration in the noise-dim ablation requires $\sim$9.5\,hours on
Apple M-series hardware (MPS) or $\sim$2--3\,hours on CUDA.
The ablation results in Table~\ref{tab:noise_results} require 5 full
training runs.

\textbf{Population diversity.}
Both datasets consist of university students in CS1 courses.
The model may not generalize to expert programmers or non-programmers.


%%─────────────────────────────────────────────────────────────────────────────
\section{Conclusion}
%%─────────────────────────────────────────────────────────────────────────────

We have presented a systematic study of conditioning in LSTM-GANs for
keystroke timing synthesis.
The key finding is that \emph{noise-to-context ratio} in the LSTM input
is the primary determinant of whether conditioning signals influence
generation.
With an 8:1 ratio (current baseline), the model achieves excellent marginal
distribution matching (KS\,=\,0.0716) but exhibits conditioning collapse
in per-state evaluation.
Reducing the ratio to 2:1 (proposed Config~D) is expected to restore
conditioning effectiveness while maintaining competitive KS.

We contribute three new evaluation tools beyond the marginal KS test
(\texttt{per\_state\_ks.py}, \texttt{temporal\_eval.py},
\texttt{ablation\_noise\_dim.py}) and recommend their adoption as a
standard diagnostic suite for conditional time series GANs.

\bibliographystyle{ACM-Reference-Format}
\bibliography{model_refs}

\end{document}
